\documentclass[10pt,a4paper]{amsart}
\usepackage[margin=19mm]{geometry}
\usepackage{amsmath,amssymb,mathtools}
\usepackage[scr=rsfs,cal=euler]{mathalfa}
\usepackage{xfrac}
\usepackage[unicode,hidelinks,bookmarks=false]{hyperref}
\newcommand{\ie}{i.e.\ }
\newcommand{\cf}{cf.\ }
\newcommand{\resp}{respectively\ }
\newcommand{\Subsection}{\subsection}
\providecommand{\nobreakdash}{\nobreak\hbox{-}}
\setlength{\emergencystretch}{2em}
\multlinegap0pt
\usepackage{amssymb}
\usepackage{mathtools}
\usepackage[shortlabels]{enumitem}
\usepackage{xcolor}
\newtheorem{lemma}{Lemma}
\title{Hamilton decompositions of the directed 5-torus for odd modulus}
\author{SangHyun Park}
\date{Source: arXiv:2604.27140v1 (CC BY 4.0)}
\begin{document}
\maketitle

\noindent\emph{Source and licence.} Hamilton decompositions of the directed 5-torus for odd modulus. Version arXiv:2604.27140v1 (CC BY 4.0), distributed by arXiv under the Creative Commons Attribution 4.0 International licence, http://creativecommons.org/licenses/by/4.0/, as recorded in the arXiv licence metadata for this exact version and shown on the abstract page https://arxiv.org/abs/2604.27140v1. Full licence notice: \texttt{licenses/arxiv-2604.27140-CC-BY-4.0.txt}.\par\smallskip

\noindent\textit{Editorial scope.} Counted unit: the article's lemma lem:longwrapclosing (source0.tex lines 667--673). The article's complete original proof of it is delivered with it (source0.tex lines 675--787, one \texttt{proof} environment of 113 lines). No mathematics is written, repaired or re-derived: the statement and the proof are the article's own lines, in the article's own notation, and no retained line is substituted or re-flowed.  No further source-local context is needed: every cross-reference of every retained line resolves inside the two delivered spans.  Nothing was omitted from any retained span, so the package carries no omission record.  No retained line contains a citation, so the wrapper carries no bibliography and no bibliography entry.  No retained line contains a figure environment, an image-inclusion command or drawing code, so no image is delivered and the package declares no asset dependency.  Editorial formatting only, in addition: an AMS \texttt{amsart} preamble that restates the article's own declarations and macros used by the retained lines, namely the environments \texttt{lemma} on separate counters, together with the standard packages those lines need.  The retained environments print in this document as Lemma 1 in order of appearance, whereas the article prints the same items with the numbering of its own full document.  The complete native source of the article as distributed in the arXiv e-print for this exact version is bundled unchanged as \texttt{sources/199448/source0.tex}.
\begin{lemma}[Long wrap skeleton]\label{lem:longwrapclosing}
The remaining point satisfies
\[
\Phi(0,m-1)=(1,0),\qquad
\ell(0,m-1)=m^3-(m-1)(m-2).
\]
\end{lemma}

\begin{proof}
The orbit from \(x_*=w(0,m-1)\) first passes through the all-zero point and then enters the family
\[
E(u,v)=(u,v,0,0,-u-v),
\]
with \(u\ne0\) and \(u+v\ne0\).  Inside this family three transitions appear: a quick seam step when
\(v+1=0\), a length-\((m-1)\) generic step when neither seam is hit, and a long seam step of length
\(3m-2\) when \(v+1=-u\).  The formulas \textup{(\(LW_1\))}--\textup{(\(LW_3\))} below give these three moves
explicitly.

At every such point the table gives \(p=1\), since the zero-set
contains \(\{2,3\}\), does not contain \(4\), and does not contain \(0\).  The long wrap starts from
\[
x_*=w(0,m-1)=(0,0,-1,0,1).
\]
The initial direction word is
\[
2,\quad 4^{m-1},\quad 0,\quad 4^{m-1}.
\]
Indeed, the first step is the section step \(p=2\).  Then \(w_1=w_2=0\) and \(w_3\ne0\), so the table
gives \(p=4\) until the all-zero point is reached after \(m-1\) steps.  At the all-zero point the table
gives \(p=0\), and then the same \(p=4\) run of length \(m-1\) ends at
\[
G^{2m}x_*=E(1,0).
\]
This initial segment has no intermediate point of \(\Sigma\): whenever \(w_3=0\) before the endpoint,
one also has \(w_4=0\).

It remains to follow the orbit from the points \(E(u,v)\).  Put \(V=v+1\).  The first step from
\(E(u,v)\) has selector \(p=1\) and gives
\[
(u-3,V,0,1,1-u-v)=(u-3,V,0,1,2-u-V).
\]
If \(V=0\), then \(w_1=0\) and \(w_3\ne0\) until the next boundary, so the next \(m-1\) selectors are
all \(4\).  Therefore
\[
G^mE(u,m-1)=E(u,0). \tag{$LW_1$}
\]

Assume now that \(V\ne0\).  Until \(w_3\) next becomes zero, the coordinate \(w_1=V\) is fixed.  Hence
the only possible nonzero selector during this flight is \(p=3\), and it occurs exactly when
\(w_4=0\); all other steps have selector \(p=0\).  Since both \(p=0\) and \(p=3\) increase \(w_4\) by
one, after the \(p=1\) step the successive values of \(w_4\) are
\[
2-u-V+j\qquad (j=0,1,2,\ldots).
\]
Thus the \(p=3\) events occur at the values \(j\equiv u+V-2\pmod m\).

If \(V\ne -u\), then the first such \(j\) lies in \(\{0,
\ldots,m-3\}\); the excluded value \(j=m-1\) would mean \(u+v=0\), which is not an \(E\)-state, and
\(j=m-2\) is exactly the case \(V=-u\).  Hence there is exactly one \(p=3\) event among the next
\(m-2\) steps.  At that time
\[
w_3=1+(m-2)+1=0,
\qquad
w_4=(2-u-V)+(m-2)=-u-V\ne0,
\]
and the first coordinate is
\[
u-3-2(m-2)-1=u.
\]
Before this endpoint, after \(k<m-2\) steps of the flight, the coordinate \(w_3\) is represented by
\(1+k+\varepsilon\) with \(\varepsilon\in\{0,1\}\); hence it lies in \(\{1,\ldots,m-1\}\) and is not zero.
Therefore
\[
G^{m-1}E(u,v)=E(u,v+1)\qquad (v+1\ne0,\; v+1\ne -u). \tag{$LW_2$}
\]

Finally suppose \(V=-u\), i.e. \(v=-u-1\).  Then \(w_4\) starts at \(2\), so the first two \(p=3\)
events occur at \(j=m-2\) and \(j=2m-2\).  The only intermediate vanishing of \(w_3\) occurs at
\(j=2m-2\), where \(w_4=0\); hence that point is neither in \(\Sigma\) nor in the \(E\)-family, and the
selector there is again \(p=3\).  After \(3m-3\) steps following the initial \(p=1\) step, there have
been exactly two \(p=3\) events, and
\[
w_3=1+(3m-3)+2=0,
\qquad
w_4=2+(3m-3)=-1,
\]
while
\[
w_0=u-3-2(3m-3)-2=u+1.
\]
Thus
\[
G^{3m-2}E(u,-u-1)=
\begin{cases}
E(u+1,-u),&u\ne m-1,\\
w(1,0),&u=m-1.
\end{cases} \tag{$LW_3$}
\]
In each of \textup{(\(LW_1\))}--\textup{(\(LW_3\))}, the displayed endpoint is the first time in the
segment with \(w_3=0\) and \(w_4\ne0\); in the long case the earlier time with \(w_3=0\) has
\(w_4=0\).

Starting from \(E(1,0)\), the rules \textup{(\(LW_1\))}--\textup{(\(LW_3\))} visit, for each fixed
\(u=1,2,\ldots,m-1\), the \(m-1\) values
\[
v=1-u,\; 2-u,\; \ldots,\; -u-1
\]
cyclically, omitting only \(v=-u\).  The last seam, from \(E(m-1,0)\), lands at \(w(1,0)\).  Hence
\(\Phi(0,m-1)=(1,0)\).

The length is the sum of the displayed segments.  The initial segment has length \(2m\).  Among the
\((m-1)^2\) points \(E(u,v)\), the case \(v+1=0\) occurs \(m-2\) times, the seam case \(v+1=-u\)
occurs \(m-1\) times, and the remaining cases occur \((m-2)^2\) times.  Therefore
\[
\ell(0,m-1)=2m+(m-2)m+(m-1)(3m-2)+(m-2)^2(m-1).
\]
A simplification gives
\[
\ell(0,m-1)=m^3-m^2+3m-2=m^3-(m-1)(m-2).
\]
\end{proof}

\end{document}
